\section{From Alerts to \trupoc{}s}
\label{sec:alert-ranking}

When disclosing a vulnerability to vendors, they often require a PoC to demonstrate that the vulnerability exists and is triggerable.
As such, \mango focuses on finding \trupoc{}s.
We design a filtering system in \mango that takes vulnerability alerts it finds and generates \trupoc{}s.
These vulnerabilities have clearly defined paths through single or multiple programs stemming from user-controlled sources to sinks.

During static analysis described in Section~\ref{sec:rda}, we build a dependency graph of data flows from all sources to sinks.
We use this dependency graph to discover and categorize possible interesting sources of input.

We categorize potential external data sources into five categories:
\textit{file\_ops} consists of file operations, such as \texttt{fgets}, \texttt{fread}, etc.
\textit{network\_ops} covers \texttt{socket}.
\textit{env\_ops} contains all references to \texttt{frontend\_params}, \texttt{getenv} and \texttt{nvram\_get} style operations, which retrieve stored data.
\textit{argv} denotes paths that appear to originate from a program's command line.
Finally, \textit{unknown} is a category used for all data from any sources not listed above.

In the context of firmware, it is easy to identify sources for \textit{env\_ops}.
If the \textit{env\_op} of an alert has a known location that sets its value or it is determined to be a user-controllable \texttt{frontend\_param}, then \mango will report the alert as a \trupoc.
% Although a \trupoc can have many data sources, \mango considers the highest-valued data source during filtering.
